% ---------------------------------------------------------------------
% 第 4 章 几何一致误差体系
% ---------------------------------------------------------------------
\section{几何一致误差体系}

\subsection{误差对象的阶数}

计算力矩律的反馈项为 $-K_de_\xi-A^\top K_pe_z$，只需位姿与速度误差；期望加速度用于前馈，动力学失配归入 $\dot e_\xi$ 中的扰动。因此正运动学与期望轨迹用 TODQ 建模，误差体系定义在 HDQ 上，无需测量加速度误差。

\subsection{定理 1：误差的 HDQ 表示}

实测链取 HDQ 表示 $\breve x=\hat{\underline x}+\varepsilon^*\dot{\hat{\underline x}}$；期望链取 $\breve x_d=\hat{\underline x}_d+\varepsilon^*\dot{\hat{\underline x}}_d$。

\begin{theorem}[误差的 HDQ 提升]\label{thm:1}
本文采用右误差 $\tilde x=x x_d^*$，并使用空间表示的 twist。定义 HDQ 误差元素
\begin{equation}
\breve{\tilde x}\triangleq \breve x\cdot\bigl(\breve x_d\bigr)^{*} .
\tag{4.1}\label{eq:4.1}
\end{equation}
则
\begin{equation}
\breve{\tilde x}=\tilde x+\varepsilon^*\dot{\tilde x},\qquad
\tilde x=\hat{\underline x}\hat{\underline x}_d^{\,*},\quad
\dot{\tilde x}=\dot{\hat{\underline x}}\hat{\underline x}_d^{\,*}+\hat{\underline x}\dot{\hat{\underline x}}_d^{\,*},
\tag{4.2}\label{eq:4.2}
\end{equation}
即一次 HDQ 乘法（3 次 DQ 乘）同时给出位姿误差 $\tilde x$ 与其导数；0 阶项正是 \cite{P2} 的 $\tilde x$。进一步定义\textbf{几何一致 twist 误差}
\begin{equation}
\tilde{\boldsymbol\xi}\triangleq 2\dot{\tilde x}\tilde x^{*},\qquad
e_\xi\triangleq\mathrm{vec}_6\tilde{\boldsymbol\xi},
\tag{4.3}\label{eq:4.3}
\end{equation}
则：(i) $\tilde{\boldsymbol\xi}$ 是纯 DQ，且对 unwinding 翻转 $\tilde x\to-\tilde x$ 不变；(ii) 误差满足与原系统同型的左乘运动学 $\dot{\tilde x}=\tfrac12\tilde{\boldsymbol\xi}\tilde x$；(iii) 由误差定义和运动学关系有
\begin{equation}
\tilde{\boldsymbol\xi}=\boldsymbol\xi-\mathrm{Ad}_{\tilde x}\boldsymbol\xi_d ,
\tag{4.4}\label{eq:4.4}
\end{equation}
即“当前 twist 减去\textbf{搬运到当前位姿处}的期望 twist”。
\end{theorem}

\begin{proof}
式 \eqref{eq:4.2} 由 HDQ 乘法与 Leibniz 法则直接得到。(i) 对 $\tilde x\tilde x^*=1$ 求导可得 $\tilde{\boldsymbol\xi}+\tilde{\boldsymbol\xi}^*=0$，故其为纯 DQ；同时翻转 $\tilde x,\dot{\tilde x}$ 不改变 $2\dot{\tilde x}\tilde x^*$。(ii) 式 \eqref{eq:4.3} 右乘 $\tilde x$ 即得。(iii) 由 $\dot{\hat{\underline x}}_d^{\,*}=-\tfrac12\hat{\underline x}_d^{\,*}\boldsymbol\xi_d$ 得 $\dot{\tilde x}=\tfrac12\boldsymbol\xi\tilde x-\tfrac12\tilde x\boldsymbol\xi_d$，右乘 $2\tilde x^*$ 即 \eqref{eq:4.4}。
\end{proof}

\subsection{定理 2：输出误差运动学}

沿用 \cite{P2} 输出 $e_z=[\mathcal O;\mathcal T]\in\mathbb R^6$，$\mathcal O=-\mathrm{Im}\,\tilde r$，$\mathcal T=\tilde p$（$\tilde r,\tilde p$ 为 $\tilde x$ 的旋转/归一化平移分量，$\tilde r=\tilde\eta+\tilde\mu$）。

\begin{theorem}[输出误差运动学闭式]\label{thm:2}
记 $\tilde{\boldsymbol\xi}=\tilde\omega+\varepsilon\tilde v$，则
\begin{equation}
\dot e_z=A(\tilde x)\,e_\xi,
\qquad
A(\tilde x)=
\begin{bmatrix}
-\tfrac12\bigl(\tilde\eta I_3+[\mathcal O]_\times\bigr) & 0_3\\[2pt]
-[\mathcal T]_\times & I_3
\end{bmatrix},
\tag{4.5}\label{eq:4.5}
\end{equation}
且 $\tilde x\to1$ 时 $A\to A_0=\mathrm{diag}(-\tfrac12I_3,I_3)$，$\sigma_{\min}(A_0)=\tfrac12$。
\end{theorem}

\begin{proof}
旋转分量：由 $\dot{\tilde r}=\tfrac12\tilde\omega\tilde r$（定理 1(ii) 的四元数部），$\dot{\mathcal O}=-\mathrm{Im}(\tfrac12\tilde\omega\tilde r)=-\tfrac12(\tilde\eta I_3+[\mathcal O]_\times)\tilde\omega$（用 $\tilde\omega\times\tilde\mu=[\mathcal O]_\times\tilde\omega$）。平移分量：由 twist 约定 $\tilde v=\dot{\tilde p}+\tilde p\times\tilde\omega$ 得 $\dot{\mathcal T}=\tilde v-[\mathcal T]_\times\tilde\omega$。
\end{proof}

定理 1、2 给出二阶误差系统：$\dot e_z=Ae_\xi$，$\dot e_\xi$ 由控制律与扰动决定（§5），无需另设加速度误差状态。
